\chapter{范围}
本文件规定了 FASTQ 原始数据的描述对象、描述内容、字段要求、数据包组织方式、描述文件要求以及质量与一致性检查要求。

本文件适用于生物样本，尤其是农作物种质资源相关 FASTQ 原始数据的采集、交付、平台入库、共享交换和长期保存。

本文件不规定 FASTQ 文件本体的序列记录语法和质量值字符编码细节。FASTQ 文件本体格式应符合有关高通量测序数据序列格式标准的规定。

\chapter{规范性引用文件}
下列文件中的内容通过文中的规范性引用而构成本文件必不可少的条款。凡是注日期的引用文件，仅该日期对应的版本适用于本文件；凡是不注日期的引用文件，其最新版本（包括所有的修改单）适用于本文件。

GB/T 1.1—2020\quad 标准化工作导则\ 第1部分：标准化文件的结构和起草规则

GB/T 35890—2018\quad 高通量测序数据序列格式规范

\chapter{术语和定义}
GB/T 35890—2018 界定的以及下列术语和定义适用于本文件。

\section{FASTQ 原始数据}
由高通量测序平台产生、以 FASTQ 形式保存的原始序列数据。

\section{FASTQ 原始数据描述}
对 FASTQ 原始数据相关的样本、测序、文件、校验及附属信息进行结构化表达的过程。

\section{描述文件}
承载 FASTQ 原始数据描述信息的结构化文件。

\section{数据包}
以样本或批次为单位组织的 FASTQ 文件及其描述文件、校验文件和附属文件的集合。

\section{样品编号}
用于样本目录组织和交付定位的编号。

\section{样本编号（分析用）}
用于后续序列比对、变异检测、基因型构建及其他分析处理的样本标识符。

\section{文库布局}
用于表示测序片段读取配置方式的信息。

\chapter{缩略语}
下列缩略语适用于本文件。

\begin{itemize}
    \item FASTQ：序列与质量信息格式；
    \item MD5：消息摘要算法；
    \item PE：Paired-End，双端测序；
    \item SE：Single-End，单端测序；
    \item VCF：Variant Call Format，变异检测文件格式；
    \item gVCF：genomic Variant Call Format，基因组变异调用文件格式。
\end{itemize}

\chapter{描述原则}
\section{完整性}
描述内容应覆盖样本追溯信息、测序与文库信息、FASTQ 文件信息、校验信息以及必要的附属信息。

\section{一致性}
描述文件中的字段值应与实际交付文件、目录结构和文件命名保持一致。

\section{可追溯性}
描述信息应能够追溯至样本、资源、批次和责任主体。

\section{可交换性}
字段命名、值域和目录组织方式应统一，便于平台交换、程序解析和长期维护。

\section{可扩展性}
在满足本文件规定的最小字段集合基础上，可根据业务需要扩展字段，但不应改变本文件规定字段的基本含义。

\chapter{描述对象与层次}
\section{总体要求}
FASTQ 原始数据描述应包括样本层、实验层、文件层和附属层。

\section{样本层}
样本层用于描述作物名称、样品编号、资源编号、种质名称、统一编号、种质类型、拉丁文名等样本追溯信息。

\section{实验层}
实验层用于描述测序平台、文库构建、测序技术、实验材料来源类型、方法选择、文库布局、读长和质量编码等信息。

\section{文件层}
文件层用于描述 FASTQ 文件名、文件格式、文件后缀、文件大小、MD5 校验码以及可选的完整性检查结果。

\section{附属层}
附属层用于描述参考序列文件、变异检测文件、备注文件以及提交单位、联系人、邮箱和提交日期等信息。

\section{描述单元}
FASTQ 原始数据描述宜采用一个样本一条记录的方式。

对于双端测序样本，应在同一条记录中分别描述 R1 与 R2 文件。

对于单端测序样本，仅描述单个 FASTQ 文件。

\chapter{描述内容要求}
\section{样本追溯信息}
\subsection{必要字段}
样本追溯信息应至少包括下列字段：

作物名称；

样品编号；

资源编号；

种质名称；

是否为全国统一编目资源；

种质类型。

\subsection{条件字段}
当“是否为全国统一编目资源”为“是”时，“原全国统一编号”和“全国统一编号”至少应填写其中一项。

\subsection{可选字段}
样本追溯信息可包括拉丁文名、鉴定时间、数据负责人和负责人联系方式等字段。

\section{实验与文库信息}
\subsection{必要字段}
实验与文库信息应至少包括下列字段：

测序平台；

文库构建；

测序技术；

文库的实验材料来源类型；

方法选择；

测序片段读取配置相关信息；

质量编码；

F 端 reads 读长。

\subsection{条件字段}
当测序片段读取配置相关信息为 PAIRED 时，应同时提供：

R 端 reads 读长；

实体数据2名称（R2）；

MD5码2。

\subsection{取值要求}
文库布局应与实际测序方式一致。

双端测序应使用 PAIRED 表达；单端测序应使用 SINGLE 或与既有系统兼容的等效表达。

\section{FASTQ 文件信息}
\subsection{必要字段}
每个样本应至少描述下列文件信息：

实体数据1名称；

MD5码1。

\subsection{双端要求}
对于双端测序样本，还应描述：

实体数据2名称；

MD5码2。

\subsection{推荐字段}
条件允许时，宜补充下列信息：

文件后缀；

文件大小；

gzip 完整性检查结果；

FASTQ 结构检查结果；

序列与质量长度一致性检查结果；

reads 数量。

\section{参考序列信息}
\subsection{条件描述}
当同时上传参考序列实体数据时，应描述：

是否同时上传参考序列实体数据；

参考序列实体数据地址；

参考序列 MD5 码。

\subsection{推荐描述}
宜补充参考序列版本和参考序列链接。

\section{变异检测文件信息}
\subsection{条件描述}
当同时上传变异检测文件时，应描述：

是否上传变异检测文件；

变异检测文件类型；

变异检测文件名；

变异检测文件 MD5。

\subsection{类型要求}
变异检测文件类型宜区分为：

单样本结果；

多样本合并结果。

\section{管理信息}
描述文件宜包括：

项目/批次编号；

提交单位；

联系人；

邮箱；

提交日期；

文库类型；

参考基因组版本；

批次备注。

\chapter{数据包组织与命名}
\section{目录组织}
\subsection{样本目录}
数据包根目录下宜按样品编号建立子目录。样本目录内放置该样本对应的 FASTQ 文件和备注文件。

\subsection{参考序列目录}
当提交参考序列文件时，宜在根目录建立 Ref 目录集中存放。

\subsection{变异检测目录}
当提交多样本合并的变异检测文件时，宜在根目录建立 VCF 目录集中存放。

\section{文件命名}
\subsection{基本要求}
FASTQ 文件名应能够区分样本及 reads 端。

\subsection{双端标识}
双端测序文件宜使用 R1、R2 或 mate1、mate2 等明确标识。

\subsection{命名字符}
文件名宜仅使用字母、数字、下划线和短横线，不宜使用空格及其他特殊字符。

\subsection{扩展名}
压缩文件扩展名宜为 .fastq.gz 或 .fq.gz。

\chapter{描述文件要求}
\section{文件形式}
描述文件可采用 XLSX、CSV、TSV 或其他结构化电子形式。

\section{字段规范}
描述文件应使用统一字段名、统一值域和统一日期格式。

\section{值域控制}
当字段采用受控词表时，应保持与预定义值域一致，不应随意改写。

\section{对应关系}
描述文件中的文件名、路径、校验码应与实际文件逐项对应。

\section{记录组织}
描述文件中的一条记录应对应一个样本。对于双端测序样本，R1 与 R2 文件应在同一条记录中分别描述。

\chapter{质量与一致性检查}
\section{基本检查项目}
描述文件与数据包之间至少应进行下列一致性检查：

a) 目录一致性：样品编号与样本目录名称一致；

b) 文件一致性：描述文件中的文件名与实际交付文件一致；

c) MD5 一致性：描述文件中的 MD5 与实际文件校验结果一致；

d) 配对一致性：Layout 为 PAIRED 时，同一样本同时存在 R1 与 R2 文件；

e) 格式合法性：FASTQ 文件结构符合 GB/T 35890—2018。

\section{推荐检查内容}
条件允许时，宜对下列内容进行检查：

gzip 完整性；

FASTQ 结构合法性；

序列与质量长度一致性；

reads 数量统计。

\section{检查结果记录}
检查结果可写入描述文件，也可形成单独的质控记录文件与数据包一并提交。

\chapter{描述文件与数据包的对应关系}
\section{样本对应关系}
描述文件中的每一条记录应能唯一对应到一个样本目录。

\section{文件对应关系}
记录中的文件名、MD5、附属文件信息应能与实际数据包中的文件逐项对应。

\section{批次覆盖关系}
当一个批次内存在多个样本时，描述文件应完整覆盖该批次所有样本。
